% ======================================================================
% paper_style_preamble.tex — drop-in style fragment (Carrasco 2026 look)
% ======================================================================
% Anchor: carrasco2026.pdf (repo root). Colors below were extracted from
% that PDF's content streams, not eyeballed:
%   citations / URLs   RGB(162, 99, 49)  = #A26331 (sienna)
%   internal links     RGB( 48, 73, 89)  = #304959 (dark slate)
%   figure series      #3C5B6F (slate navy) / #BA8F6B (muted copper),
%                      darker accents #1A476F / #B36633; gridlines 50% grey
% The matplotlib mirror of the figure palette is paper_plot_style.py.
%
% HOW TO USE WITH Draft_august.tex
%   1. DELETE (or comment out) \usepackage{lmodern} in the main preamble.
%      Palatino (newpx) and Latin Modern fight over the text font; the
%      last one loaded wins and the metrics disagree. This fragment loads
%      newpxtext/newpxmath (fallback: mathpazo).
%   2. \input{style/paper_style_preamble} AFTER the
%      existing package block. It is compatible with the rest of
%      Draft_august's preamble (amsmath/amssymb/mathtools, geometry,
%      setspace, graphicx, booktabs, caption, subcaption, float,
%      appendix, natbib, xcolor): everything here either loads a package
%      not already present or only calls \hypersetup / \captionsetup.
%   3. DELETE the existing \hypersetup{colorlinks=true, linkcolor=blue,
%      citecolor=blue, urlcolor=blue} line — this fragment replaces it.
%      (If hyperref is not yet loaded, this fragment loads it.)
% ======================================================================

% ---- Typography: Palatino text + matching math (Carrasco uses the
%      Palatino family; newpx is the maintained implementation) ---------
\IfFileExists{newpxtext.sty}{%
  \usepackage{newpxtext}%
  \usepackage{newpxmath}%
}{%
  \usepackage{mathpazo}% fallback: Palatino text + math
}
\linespread{1.18} % Palatino needs more leading; approximates the anchor's
                  % openness without double-spacing (setspace, if loaded,
                  % may override this — call \onehalfspacing after if desired)
\setlength{\emergencystretch}{1.5em} % Palatino sets wider than CM/LM;
                  % this absorbs the overfull lines the font swap creates

% ---- Colors (extracted from carrasco2026.pdf) ------------------------
\usepackage{xcolor} % no-op if already loaded
\definecolor{citesienna}{RGB}{162, 99, 49}   % citations, URLs
\definecolor{linkslate}{RGB}{48, 73, 89}     % internal refs (sections, eqs, tables)
\definecolor{sernavy}{RGB}{26, 71, 111}      % figure series 1 (dark navy)
\definecolor{sercopper}{RGB}{179, 102, 51}   % figure series 2 (burnt sienna)
\definecolor{sernavyline}{RGB}{60, 91, 111}  % muted line variant of series 1
\definecolor{sercopperline}{RGB}{186, 143, 107} % muted line variant of series 2

% ---- Hyperref: colored author-year citations, slate internal links ---
\makeatletter
\@ifpackageloaded{hyperref}{}{\usepackage{hyperref}}
\makeatother
\hypersetup{
  colorlinks = true,
  citecolor  = citesienna,
  urlcolor   = citesienna,
  linkcolor  = linkslate,
  filecolor  = linkslate,
}

% ---- Captions: "Table 1: Baseline calibration" above tables/figures,
%      slightly smaller than body, plain weight (anchor style) ---------
\makeatletter
\@ifpackageloaded{caption}{}{\usepackage{caption}}
\makeatother
\captionsetup{labelsep=colon, font=small, labelfont=bf, skip=6pt}
% Anchor convention: captions sit ABOVE tables and ABOVE figures.
% (Position is set where the float is written, not here.)

% ---- Tables: booktabs + the Carrasco Table-1 idiom -------------------
\makeatletter
\@ifpackageloaded{booktabs}{}{\usepackage{booktabs}}
\makeatother

% Small-caps column heads: \thead{Description} -> DESCRIPTION in
% spaced small caps, the anchor's header style.
\newcommand{\thead}[1]{\multicolumn{1}{c}{\footnotesize\textsc{#1}}}
\newcommand{\theadl}[1]{\multicolumn{1}{l}{\footnotesize\textsc{#1}}}
\newcommand{\theadr}[1]{\multicolumn{1}{r}{\footnotesize\textsc{#1}}}

% Bold full-width group-separator row: \tgroup{4}{Demographics and Healthcare}
% (first argument = number of columns in the tabular).
\newcommand{\tgroup}[2]{\multicolumn{#1}{l}{\textbf{#2}}\\}

% Under-table notes minipage: \tabnote{0.92}{...} — bold run-in "Notes."
% then footnotesize flowing text, width as a fraction of \textwidth.
\newcommand{\tabnote}[2]{%
  \vspace{0.5em}\par
  \begin{minipage}{#1\textwidth}
  \footnotesize\emph{Notes.} #2
  \end{minipage}}

% Under-figure note, anchor style ("Note." then series names in their
% colors): \fignote{0.9}{Expenditure ... \serone{Top 10 percent} and
% \sertwo{Bottom 90 percent}.}
\newcommand{\fignote}[2]{%
  \vspace{0.4em}\par
  \begin{minipage}{#1\textwidth}
  \footnotesize\textbf{Note.} #2
  \end{minipage}}
\newcommand{\serone}[1]{\textcolor{sernavy}{#1}}
\newcommand{\sertwo}[1]{\textcolor{sercopper}{#1}}

% ----------------------------------------------------------------------
% TEMPLATE — the Carrasco Table-1 skeleton (baseline-calibration table).
% Copy, uncomment, fill. Structure: caption ABOVE; four columns
% (symbol | description | value | target/source); small-caps heads;
% bold group rows separating parameter families; sources as \citep{}
% (they render sienna); notes minipage below.
% ----------------------------------------------------------------------
% \begin{table}[H]
% \centering
% \caption{Baseline calibration}
% \label{tab:baseline_calibration}
% \begin{tabular}{@{}llcl@{}}
% \toprule
%  & \theadl{Description} & \thead{Value} & \theadr{Target/Source} \\
% \midrule
% \tgroup{4}{Demographics and Healthcare}
% $n$              & Fertility rate            & $0.03$      & Population growth $=1\%$ \\
% $\bar\psi$       & Survival rate parameters  & See text    & WPP, HRS, MEPS \\
% \tgroup{4}{Preferences}
% $\beta$          & Subjective discount rate  & Internally calibrated & Interest rate $4\%$ (p.a.) \\
% $\sigma$         & Risk aversion             & $2.0$       & \citep{SomeSource2019} \\
% \bottomrule
% \end{tabular}
% \tabnote{0.92}{Rates are annual. Internally calibrated parameters are
% estimated in Section~X; see Table~Y.}
% \end{table}
% ----------------------------------------------------------------------
